% U5-EN — English review manuscript, part 1 (Sections 1–3)
% Derivative of the reviewed Korean manuscript at ../u3_kros/ (adaptation, not literal translation).
% Body only: no preamble. Macros used: \jj, \dd (defined in the main file).
%
% ABSTRACT: With the spread of collaborative robots and humanoids, users who
% are not control specialists increasingly need to configure how a robot
% behaves in contact. The core language of contact behavior is impedance, yet
% most of the literature presents it at an expert level with derivations
% omitted, forcing newcomers to hop between references to reconstruct a
% single concept. Part 1 of this review starts from the historical problem in
% which the language of impedance grew up --- signal distortion on long
% telegraph lines --- and then develops, without skipping derivation steps and
% in a single consistent notation, the definition and differentiation rule of
% the Laplace transform, the electrical impedance of RLC circuits, and the
% mechanical impedance of the mass-spring-damper system. Every numerical
% example can be reproduced in an open educational MuJoCo simulator, and an
% extended manuscript containing all intermediate steps is provided in a
% public repository together with the simulator. Part 2 builds on this
% foundation to treat robot impedance control and the principles of safe
% parameter selection.

\section{Introduction}\label{sec:u5-intro}

Robots have left the factory floor. In homes, hospitals, and small workshops
it is becoming routine for a robot to pick up objects, wipe surfaces, and
insert parts in spaces shared with people \cite{nvidia_physical_ai_2025}.
For such contact tasks, quality and safety are not determined by positioning
accuracy alone; they depend on how the robot has been configured to relate
force to motion --- that is, on its impedance
\cite{hogan1985impedancepart1,abudakka2020variableimpedance}.
The difficulty is that this configuration is not easy even for specialists.
Combine a stiffness with a target offset carelessly and the elastic energy
stored in the virtual spring is released all at once, sending the robot
flying; misunderstand the notion of an equilibrium point and perfectly normal
behavior gets misread as a malfunction.
The author went through exactly this kind of trial and error in front of an
expensive physical robot.

There are two reasons why the existing literature offers little relief.
The first is omitted derivations.
Papers and textbooks skip intermediate steps because of page limits and
assumptions about the reader, and a newcomer ends up wandering across several
references in search of the one step that was left out.
The second is inconsistent notation.
The same physical quantity appears under different symbols in different
sources, so study time is spent matching symbols rather than understanding
concepts.

This review confronts both problems head-on.
It begins at a historical entrance --- showing that impedance is a language
that grew out of a real engineering problem, the signal distortion of
nineteenth-century long-distance telegraph lines --- then derives the Laplace
transform from its definition, and connects electrical (RLC) and mechanical
(mass-spring-damper) impedance in one consistent notation.
For reasons of space, the body of the paper carries the essential
derivations; an extended manuscript (122 pages) containing every intermediate
step, together with an educational MuJoCo simulator that reproduces all
numerical examples, is provided in a public repository.
Part 2 builds on this foundation to cover impedance control, criteria for
choosing between impedance and admittance control, and principles of safe
parameter selection grounded in failure cases.

\section{The Historical Origin of Impedance}\label{sec:u5-history}

Etymologically, the word \emph{impedance} carries the flavor of hindering
progress, but in engineering it means more than simple obstruction: it is a
quantity that expresses both how a system resists an applied flow and how it
stores energy and gives it back.
Why such a concept became necessary is clearest if we return to the problems
of long-distance telegraph and telephone lines in the late nineteenth
century.

\subsection{The long line that broke the resistance-only picture}

For a short DC circuit on a benchtop, one can pack the entire notion of
opposition to flow into a single resistance and explain a great deal with
Ohm's law $V=IR$ alone.
But a telegram is a signal of short pulses switching on and off, and a
telephone call, like a human voice, is a signal that changes continuously;
neither is DC.
On real long cables something strange happened.
Pulses sent over great distances arrived late and rounded off; voices had
their high and low components transmitted differently, so that far away
they arrived smeared and hard to understand.
The sharper the edge of a pulse, the more high-frequency content it needs,
so if the line weakens the high components more, or makes different
frequencies arrive at different times, the distinction between dots and
dashes blurs.
The DC-style intuition --- just lower the wire resistance --- could not
explain this.

The reason is that a long wire is not a simple resistance.
Along a long line there are, distributed with position, a resistance per
unit length $R'~[\Omega/\mathrm{m}]$, an inductance $L'~[\mathrm{H/m}]$, a
capacitance $C'~[\mathrm{F/m}]$, and a conductance $G'~[\mathrm{S/m}]$
representing insulation leakage.
$R'$ and $G'$ are loss elements through which energy disappears as heat and
leakage; $L'$ and $C'$ are storage elements that hold energy briefly in
magnetic and electric fields and return it.
Unlike a lumped-element model that gathers all effects at a single point, a
long line must be viewed as a distributed system in which these small
effects repeat along its length and accumulate.
Writing this idea in terms of the voltage $v(x,t)$ and current $i(x,t)$ at
position $x$ and time $t$ gives the basic form of the telegrapher's
equations,
\begin{align}
  \frac{\partial v}{\partial x} &= -R'i - L'\frac{\partial i}{\partial t}, \\
  \frac{\partial i}{\partial x} &= -G'v - C'\frac{\partial v}{\partial t}.
\end{align}
The first equation says that the voltage changes along the line because of
series resistive loss and the series inductance voltage; the second says
that the current changes because of shunt leakage current and capacitor
charging current.
There is no need to solve these equations here.
It is enough to read them as a sign that a long wire is not one lump of
resistance but an endless chain of small storage and loss elements.

\subsection{Heaviside and the distortionless condition}

Oliver Heaviside encountered this problem first-hand while working as a
telegraph operator, and by applying Maxwell's electromagnetic theory to the
communication-line problem he treated the distortion of long lines
mathematically \cite{heaviside_electrical_papers,donaghyspargo2018heaviside}.
His insight was that the goal is not to send the signal as strongly as
possible, but to make the various frequency components weaken and delay
according to similar rules.
A waveform is a sum of frequency components, so if only some components are
attenuated more or delayed more, the overall shape collapses.
For an idealized line that is uniform and linear, with $R'$, $L'$, $C'$, and
$G'$ nearly independent of frequency, the condition for distortion-free
transmission is
\begin{equation}
  \frac{R'}{L'} = \frac{G'}{C'}.
\end{equation}
$R'/L'$ measures, on the series side, how large the loss is relative to
magnetic-field storage; $G'/C'$ measures, on the shunt side, how large the
leakage is relative to electric-field storage.
Both quantities have units of $\mathrm{s}^{-1}$, so this condition is an
equation matching time scales of the same kind.
Old telephone lines had small leakage $G'$, hence a small $G'/C'$, but
lacked inductance, so $R'/L'$ was comparatively large; adding coils at
regular intervals raised the effective $L'$ and brought the two ratios close
together over the voice band.
Adding inductance --- the very element that opposes changes in current ---
actually improved the signal; this is not an increase in obstruction but a
design that balances loss against storage to reduce waveform distortion.
So the word \emph{impedance} that Heaviside used in the 1880s was not a mere
coinage: it was a working engineer's language for explaining signal
transmission by taking into account inductive and capacitive effects that
resistance alone could not capture
\cite{ethw_heaviside,ptb_impedance_history}.

\subsection{Resistance, reactance, impedance}

For DC, it is often enough to describe the difficulty of flow by a single
resistance.
For AC, however, one must look not only at how much the signal shrinks, but
also at how much the oscillations of voltage and current lead or lag each
other in time.
Picture a practical test: send a pure tone of a single frequency into one
end of a telephone line and measure, at the other end, how much smaller it
has become and how much later it oscillates.
Inductors store energy in magnetic fields and capacitors in electric
fields, and in doing so they shift the phase between voltage and current in
a frequency-dependent way; the frequency-dependent opposition created by
such storage elements is called \emph{reactance}.
Reactance is also measured in ohms ($\Omega$), but unlike resistance it does
not mean that energy is, on average, destroyed as heat.
The AC notion of opposition therefore carries the dissipative component and
the storage component together, written as
\begin{equation}
  Z = R + \jj X.
\end{equation}
The real part $R$ is the resistance that dissipates energy, the imaginary
part $X$ is the reactance that shifts phase, and $\jj$ is the imaginary unit
marking a $90^\circ$ phase difference.
In the sinusoidal steady state the impedances of the three basic elements
are
\begin{equation}
  Z_R = R, \qquad
  Z_L = \jj\omega L, \qquad
  Z_C = \frac{1}{\jj\omega C} = -\frac{\jj}{\omega C},
\end{equation}
where $\omega=2\pi f$ is the angular frequency.
The inductor demands a larger voltage for fast current changes as $\omega$
grows, while the capacitor passes high-frequency currents relatively
easily.
In a series connection the same current flows through all three elements
and the source must supply the sum of the three element voltages, so the
total impedance is the sum of the three impedances,
\begin{equation}
  Z(\jj\omega) = Z_R + Z_L + Z_C
  = R + \jj\left(\omega L - \frac{1}{\omega C}\right).
\end{equation}
The positive reactance of the inductor and the negative reactance of the
capacitor point in opposite phase directions and subtract, and at some
frequency they cancel, leaving the circuit looking as if only the
resistance remained.
This cancellation is the seed of resonance; the detailed development of
resonance itself is treated in the open extended manuscript.
In this way, an AC circuit that would otherwise demand a differential
equation can be handled with the algebraic relation $V=IZ$, an equation
shaped like Ohm's law.
For the detailed derivation of the telegraph-line analysis and the extended
discussion of phasor notation, see the open extended manuscript.

The question the telegraph story leaves behind is how the magnitude and
phase of a time-varying input change as it passes through a system --- and to
write that question down, and to compare electrical circuits and mechanical
systems in the same sentences, we first need the grammar of the Laplace
transform and of linear time-invariant systems.

\section{A Minimal Grammar of the Laplace Transform}\label{sec:u5-laplace}

To turn the question left by the telegraph line --- why and how do a
signal's magnitude and phase change? --- into something computable, two
preparations are needed.
One is an assumption about the system; the other is a transform that turns
differential equations into algebraic ones.
This section lays out that minimal grammar without skipping any derivation
steps.

Start with the assumption.
A system is \emph{linear} when the rule $\mathcal{S}$ that maps an input
$u(t)$ to an output $y(t)$ satisfies
$\mathcal{S}\{a u_1 + b u_2\} = a\,\mathcal{S}\{u_1\} + b\,\mathcal{S}\{u_2\}$
for arbitrary inputs $u_1$, $u_2$ and scalars $a$, $b$.
Double the input and the output doubles; feed in the sum of two inputs and
the result equals the sum of the individual outputs.
A system is \emph{time-invariant} when applying the same input $t_0$ later
merely delays the output by the same $t_0$ without changing its shape,
that is, $\mathcal{S}\{u(t-t_0)\} = y(t-t_0)$.
A real robot is neither perfectly linear nor time-invariant, but for small
motions near a single operating point this linear time-invariant (LTI)
assumption explains a great many phenomena well enough, and it unlocks the
powerful tools of transfer functions and frequency responses.

\subsection{From the definition to the differentiation rule}

The Laplace transform is an integral that summarizes a time function $x(t)$
as a function of a complex variable $s$:
\begin{equation}
  X(s) = \mathcal{L}\{x(t)\}
  = \int_{0}^{\infty} x(t)\,e^{-st}\,\dd t.
  \label{eq:u5-laplace-def}
\end{equation}
In words: multiply $x(t)$ by the steadily shrinking weight $e^{-st}$ and add
everything up from $t=0$ onward.
Each choice of $s$ produces one number, so the entire time waveform is
summarized on the single sheet $X(s)$.

What makes this transform useful is that differentiation turns into
multiplication by $s$ --- and this rule is not a formula to memorize but a
consequence of the definition \eqref{eq:u5-laplace-def} and the product rule
alone.
Since $s$ does not depend on $t$,
\begin{equation}
  \frac{\dd}{\dd t}\left[x(t)\,e^{-st}\right]
  = \dot{x}(t)\,e^{-st} - s\,x(t)\,e^{-st}.
\end{equation}
Integrate both sides from $t=0$ to $\infty$.
The left side integrates a derivative, so it equals the end value minus the
start value: the value of $x(t)e^{-st}$ as $t\to\infty$ minus its value
$x(0)$ at $t=0$.
For values of $s$ at which the response does not blow up, so that
$x(t)e^{-st}\to 0$ as $t\to\infty$, the left side is $0 - x(0) = -x(0)$.
The two integrals on the right side are, by the definition,
$\mathcal{L}\{\dot{x}\}$ and $-sX(s)$ respectively, so
$-x(0) = \mathcal{L}\{\dot{x}\} - sX(s)$, and rearranging gives
\begin{equation}
  \mathcal{L}\{\dot{x}(t)\} = sX(s) - x(0).
  \label{eq:u5-deriv-rule}
\end{equation}
With the initial condition set to zero, $\dot{x}$ can be handled as
$sX(s)$.

Let us run a simple check.
If $x(t)=1$ (a constant), the definition gives
$X(s)=\int_0^\infty e^{-st}\dd t
=\left[-e^{-st}/s\right]_0^\infty
=0-\left(-\tfrac{1}{s}\right)=\tfrac{1}{s}$.
Since $\dot{x}=0$, the left side of rule \eqref{eq:u5-deriv-rule} is
$\mathcal{L}\{0\}=0$, and the right side is
$sX(s)-x(0)=s\cdot(1/s)-1=0$ as well --- they agree.

For a twice-differentiated term, apply the same rule twice.
Treating $\dot{x}$ as a new function and applying
\eqref{eq:u5-deriv-rule} again,
\begin{equation}
  \mathcal{L}\{\ddot{x}\}
  = s\,\mathcal{L}\{\dot{x}\} - \dot{x}(0)
  = s^2 X(s) - s\,x(0) - \dot{x}(0),
\end{equation}
and with zero initial conditions $\ddot{x}$ can be handled as $s^2 X(s)$.

\subsection{From differential equation to transfer function}

Now take the equation of motion of the mass-spring-damper,
$m\ddot{x}+b\dot{x}+kx=f$, and transform it term by term under zero initial
conditions:
$m\ddot{x}\to m s^2 X(s)$, $b\dot{x}\to b s X(s)$, $kx\to kX(s)$, and
$f\to F(s)$, so that
\begin{equation}
  (m s^2 + b s + k)\,X(s) = F(s).
\end{equation}
A differential ledger that had to be updated continuously in time has become
a single page of algebra.
Arranging it as a ratio of output to input gives the \emph{transfer
function},
\begin{equation}
  H(s) = \frac{X(s)}{F(s)} = \frac{1}{m s^2 + b s + k}.
  \label{eq:u5-msd-tf}
\end{equation}
In the denominator, the $ms^2$ term carries the influence of inertia, the
$bs$ term that of damping, and the $k$ term that of the spring's restoring
action.
A transfer function is an input-output relation with the initially stored
energy set to zero (zero-state), and it is worth remembering that the same
physical device presents different transfer functions depending on what is
chosen as input and output.

Here $s$ is not a mere letter of the alphabet but a complex variable that
encodes the character of the time response, usually thought of as
$s=\sigma+\jj\omega$.
$\sigma$ relates to how fast the response grows or decays, and $\omega$ to
how fast it oscillates.
This does not mean that physical inputs are literally complex numbers; it is
a representation that lets decay and oscillation be computed together
inside one symbol.

\subsection{Reading the frequency response}

Feed a stable LTI system the sinusoidal input $u(t)=A\cos(\omega t)$, and
the steady-state output has the same frequency, with only its magnitude and
phase changed:
\begin{equation}
  y_{\mathrm{ss}}(t)
  = A\,|H(\jj\omega)|\cos\!\left(\omega t + \angle H(\jj\omega)\right).
\end{equation}
$|H(\jj\omega)|$ and $\angle H(\jj\omega)$ are the magnitude and angle of
the complex number $H(\jj\omega)$.
Plot a complex number $a+\jj b$ as the point $(a,b)$ in the plane: its
magnitude is the distance from the origin, $\sqrt{a^2+b^2}$, and its angle
is the rotation measured from the positive real axis.
For example, if at some frequency $H(\jj\omega)=1-\jj$, then the magnitude
is $\sqrt{1^2+(-1)^2}=\sqrt{2}$ and the angle is $-45^\circ$, so the output
amplitude is about 1.41 times the input and the phase lags by $45^\circ$.
In other words, the single object $H(\jj\omega)$, obtained by substituting
$s=\jj\omega$, tells us at each frequency how much the system amplifies or
attenuates and how late it responds.
For the relation between poles and response shape, convolution, the
final-value theorem, and other material abbreviated here, see the extended
manuscript.

We are now ready to apply this grammar to its first physical object: in the
next section we transform the equations of the RLC circuit --- resistor,
inductor, capacitor --- by the same procedure and see how electrical
impedance is read from the result.
